\daytitle{19}{恒定电流与基本电路}

\begin{learningbox}{本节导读}
前面几节从静电场、电势差和电容器认识了“电荷如何受力、如何储能”。本节转向持续流动的电荷：从电荷的定向移动理解电流、电阻、欧姆定律，再到串联和并联电路的分析方法。这样，电势差不再只是两点能量差，也成为推动电荷在闭合电路中持续流动的原因。
\end{learningbox}

\textbf{本节读题路线：} 先标出等势节点，再把能直接化简的部分化成串并联结构；随后判断各支路电流和各元件电压怎样分配。功率题先明确是某个元件的端口功率，还是纯电阻发热功率。遇到复杂网络或非线性元件，先写节点电流、回路电压或读伏安图像，再判断 $R=U/I$ 是否能作为常量使用。

\section*{一、学习目标}

\begin{enumerate}
  \item 能用电荷定向移动解释电流 $I = \dfrac{\Delta Q}{\Delta t}$ 和 $I=n\lvert q\rvert Sv_{\mathrm d}$，并使用 $R = \dfrac{U}{I}$、$I = \dfrac{U}{R}$。
  \item 能用电阻定律 $R = \rho\dfrac{l}{S}$ 判断材料、长度和横截面积对电阻的影响。
  \item 会分析串联和并联电路中 $I$、$U$、$R$ 的分配关系。
  \item 能用电功 $W = IUt$ 和电功率 $P = IU$ 描述电能转化及其快慢。
  \item 能先识别串并联结构，再进行电路化简计算。
  \item 能从伏安特性图判断二极管等非线性元件不能直接套用欧姆定律。
  \item 能说明热敏、光敏等传感器如何通过电阻变化输出电压信号。
\end{enumerate}
\section*{二、初高中衔接}

初中已经学过这些基本知识——电流、电压、电阻、欧姆定律、串联并联。高中在以下方面深化：
\begin{enumerate}
  \item \textbf{电流的微观解释：} $I = n\lvert q\rvert Sv_{\mathrm d}$（$n$ 为载流子浓度，$\lvert q\rvert$ 为单个载流子的电荷量大小，$v_{\mathrm d}$ 为定向移动速率）。
  \item \textbf{电阻定律：} $R = \rho \dfrac{l}{S}$——电阻与材料、长度、横截面积的定量关系。
  \item \textbf{电功率的精确理解：} 计算某个元件的端口功率时用 $P = UI$；在纯电阻元件中，还可写成 $P = I^2R$ 或 $P = \dfrac{U^2}{R}$。
  \item \textbf{电路化简和复杂计算：} 多电阻串并联的等效电阻计算。
\end{enumerate}

\textbf{科学发展视角：} 真实电子设备从金属导线发展到半导体芯片，但基本电路思想仍然重要。传感器把温度、光照、压力等物理量转化为电信号，后续电路再进行放大、计算和控制。本节的串并联、功率和非纯电阻电路，是理解现代电子系统的入口。

\section*{三、核心概念}

\begin{conceptbox}{电流与电阻}

\textbf{电流：} 单位时间内通过导体横截面的电荷量。
\[
I = \frac{\Delta Q}{\Delta t}
\]
单位：安培（$\mathrm{A}$）。方向：正电荷定向移动的方向（与电子移动方向相反）。

\begin{center}
% TikZ 图：电流微观表达式的推导
% Day 19：恒定电流
\begin{tikzpicture}[scale=0.9,>=Latex]
  \draw[thick] (0,0) rectangle (6.8,1.55);
  \draw[fill=blue!8,draw=blue!60,thick] (6.8,0.775) ellipse (0.28 and 0.775);
  \node[blue,right] at (7.05,0.78) {截面 $S$};

  \fill[orange!18] (4.45,0) rectangle (6.8,1.55);
  \draw[dashed,orange!80!black,thick] (4.45,0) -- (4.45,1.55);
  \draw[<->,thick] (4.45,-0.35) -- (6.8,-0.35);
  \node[below] at (5.62,-0.35) {$v_{\mathrm d}\Delta t$};

  \foreach \x/\y in {0.55/0.35,1.25/1.05,1.9/0.55,2.55/1.18,3.2/0.35,3.85/0.96,4.75/0.42,5.25/1.10,5.85/0.72,6.35/1.25} {
    \fill[red] (\x,\y) circle (1.9pt);
    \draw[->,red] (\x+0.06,\y) -- ++(0.36,0);
  }

  \draw[->,very thick] (0.55,2.05) -- (6.25,2.05);
  \node[above] at (3.4,2.05) {约定电流方向};
  \node[red] at (3.4,1.78) {正载流子等效运动};
  \node[align=left,font=\small] at (3.2,-1.05)
    {在 $\Delta t$ 内能穿过截面 $S$ 的载流子，\\来自截面上游长度为 $v_{\mathrm d}\Delta t$ 的一段，体积 $Sv_{\mathrm d}\Delta t$。};
\end{tikzpicture}


\figcaption{电流的微观模型：载流子定向移动形成电流}\label{fig:current-microscopic}
\end{center}

若单个载流子的电荷量大小为 $\lvert q\rvert$、数密度为 $n$、导体横截面积为 $S$、定向移动速率为 $v_{\mathrm d}$，则可从“单位时间内穿过截面的载流子数”推出电流大小的微观表达式。下标 $\mathrm d$ 表示漂移（定向移动），用来区别电子远大于它的无规则热运动速率。

\textbf{推导：} 在很短时间 $\Delta t$ 内，能穿过截面 $S$ 的载流子，来自截面上游长度为 $v_{\mathrm d}\Delta t$ 的一段导体。这个小柱体体积为
\[
V=S v_{\mathrm d}\Delta t.
\]
其中载流子个数为
\[
N=nV=nSv_{\mathrm d}\Delta t.
\]
这些载流子携带的电荷量为
\[
\Delta Q=N\lvert q\rvert=n\lvert q\rvert Sv_{\mathrm d}\Delta t.
\]
代入电流定义 $I=\dfrac{\Delta Q}{\Delta t}$，得到
\[
\boxed{I=n\lvert q\rvert Sv_{\mathrm d}}.
\]
图中采用正载流子等效示意；在金属导体中，电子实际定向移动方向与约定电流方向相反。这个公式说明电流大小不仅与载流子运动快慢有关，也与载流子多少和导体粗细有关。

\textbf{欧姆定律：}
\[
\boxed{I = \frac{U}{R}}
\]

\begin{itemize}
  \item 适用条件：线性导体，且温度等条件基本不变。
  \item 不适用于二极管、气体放电管等明显非线性元件；灯丝电阻随温度变化，严格说也不是恒定电阻。
  \item 对给定导体，在温度、形状和材料状态基本不变时，$R$ 是导体本身的性质，与外加 $U$、通过的 $I$ 无关。
\end{itemize}

\textbf{实验事实与微观理解：} 实验表明，对同一段金属导体，在温度等状态基本不变、外加电压处于一定范围内时，电流与电压成正比：
\[
I\propto U.
\]
把比例系数写成 $\dfrac{1}{R}$，就得到 $I=\dfrac{U}{R}$。从微观上看，同一段导体中平均场强约为 $E=U/l$，温度基本不变时漂移速率 $v_{\mathrm d}$ 与 $E$ 近似成正比，因而 $I=n\lvert q\rvert Sv_{\mathrm d}$ 随 $U$ 成正比。这一微观图景解释了实验规律，但欧姆定律本身仍是有适用条件的实验定律，不是电阻的定义。

\textbf{电阻定律：}
\[
\boxed{R = \rho \frac{l}{S}}
\]

\begin{itemize}
  \item $\rho$：电阻率（$\Omega \cdot \mathrm{m}$），与材料和温度等状态有关。
  \item 金属的电阻率随温度升高而增大。
  \item 超导现象：某些材料温度低于临界温度时，直流电阻可降为零。
\end{itemize}

\textbf{实验规律与直观理解：} 控制材料、温度和横截面积相同，实验得到 $R\propto l$；控制材料、温度和长度相同，实验得到 $R\propto \dfrac{1}{S}$。导体越长，载流子在定向移动中受到阻碍的路程越长；横截面积越大，可同时参与导电的载流子越多。这些图景帮助理解变化方向，但不能代替控制变量实验。把材料和温度的影响归入电阻率 $\rho$，得到
\[
R=\rho\frac{l}{S}.
\]

\end{conceptbox}

\begin{conceptbox}{串联与并联电路}

\begin{center}
% TikZ 图：串联并联对比
% Day 19：直流电路
\begin{tikzpicture}[scale=0.9,>=Latex]
  % 串联
  \draw[thick] (0,0) -- (1,0);
  \draw[thick] (1,0.5) -- (1,-0.5);
  \draw[thick] (0.8,0.5) -- (1.2,0.5);
  \draw[thick] (0.8,-0.5) -- (1.2,-0.5);
  \node[above] at (1,0.8) {$R_1$};
  \draw[thick] (1,0) -- (2.5,0);
  \draw[thick] (2.5,0.5) -- (2.5,-0.5);
  \draw[thick] (2.3,0.5) -- (2.7,0.5);
  \draw[thick] (2.3,-0.5) -- (2.7,-0.5);
  \node[above] at (2.5,0.8) {$R_2$};
  \draw[thick] (2.5,0) -- (3.5,0);
  \node[below] at (1.75,-0.8) {串联：$R = R_1 + R_2$};
  
  % 并联
  \draw[thick] (5.5,0.8) -- (9.5,0.8);
  \draw[thick] (5.5,-0.1) -- (9.5,-0.1);
  \draw[thick] (7,0.8) -- (7,-0.1);
  \draw[thick] (7.5,0.8) -- (7.5,-0.1);
  \draw[thick] (7.5,0.5) -- (7.5,-0.5);
  \draw[thick] (7.3,0.5) -- (7.7,0.5);
  \draw[thick] (7.3,-0.5) -- (7.7,-0.5);
  \node[above] at (7.5,1.1) {$R_1$};
  \draw[thick] (8.5,0.8) -- (8.5,-0.1);
  \draw[thick] (8.5,0.5) -- (8.5,-0.5);
  \draw[thick] (8.3,0.5) -- (8.7,0.5);
  \draw[thick] (8.3,-0.5) -- (8.7,-0.5);
  \node[above] at (8.5,1.1) {$R_2$};
  \node[below] at (7.5,-0.8) {并联：$\dfrac{1}{R} = \dfrac{1}{R_1} + \dfrac{1}{R_2}$};
\end{tikzpicture}


\figcaption{串联（左）和并联（右）电路对比}\label{fig:series-parallel}
\end{center}

\textbf{串联电路：}
\begin{itemize}
  \item 电流处处相等：$I = I_1 = I_2$。
  \item 总电压等于各分电压之和：$U = U_1 + U_2$。
  \item 总电阻：$R = R_1 + R_2$。
  \item 电压分配：$\dfrac{U_1}{U_2} = \dfrac{R_1}{R_2}$（串联分压）。
\end{itemize}

\textbf{并联电路：}
\begin{itemize}
  \item 各支路电压相等：$U = U_1 = U_2$。
  \item 总电流等于各支路电流之和：$I = I_1 + I_2$。
  \item 总电阻：$\dfrac{1}{R} = \dfrac{1}{R_1} + \dfrac{1}{R_2}$。
  \item 电流分配：$\dfrac{I_1}{I_2} = \dfrac{R_2}{R_1}$（并联分流）。
\end{itemize}

\end{conceptbox}

\begin{conceptbox}{节点与回路：电路方程的守恒思想}

\begin{center}
% TikZ 图：节点电流守恒与回路电压守恒
% Day 19：直流电路
\begin{tikzpicture}[scale=0.92,>=Latex]
  \tikzset{
    wire/.style={thick},
    current/.style={->,very thick,red},
    volt/.style={->,thick,blue},
    res/.style={thick,fill=white},
    note/.style={font=\small}
  }

  % 左：节点电流守恒
  \begin{scope}
    \node[font=\bfseries] at (1.9,1.75) {节点电流守恒};
    \coordinate (O) at (2,0);
    \filldraw[black] (O) circle (2.2pt);

    \draw[wire] (-0.2,1.0) -- (O);
    \draw[wire] (-0.2,-1.0) -- (O);
    \draw[wire] (O) -- (4.2,0);

    \draw[current] (0.15,0.82) -- (1.25,0.27);
    \node[red,above left] at (0.50,0.88) {$I_1$};
    \draw[current] (0.15,-0.82) -- (1.25,-0.27);
    \node[red,below left] at (0.50,-0.88) {$I_2$};
    \draw[current] (2.45,0.22) -- (3.55,0.22);
    \node[red,above] at (3.00,0.24) {$I_3$};

    \node[note,align=center] at (2,-1.55) {节点不积累电荷\\[2pt]$I_1+I_2=I_3$};
  \end{scope}

  % 右：回路电压守恒
  \begin{scope}[xshift=5.6cm]
    \node[font=\bfseries] at (2.45,2.05) {回路电压守恒};

    \draw[wire] (0,-0.58) -- (0,1.15) -- (1.0,1.15);
    \draw[res] (1.0,0.85) rectangle (2.0,1.45);
    \node at (1.5,1.15) {$R_1$};
    \draw[wire] (2.0,1.15) -- (3.1,1.15);
    \draw[res] (3.1,0.85) rectangle (4.1,1.45);
    \node at (3.6,1.15) {$R_2$};
    \draw[wire] (4.1,1.15) -- (5.0,1.15) -- (5.0,-1.05) -- (0,-1.05) -- (0,-0.82);

    % 电源
    \draw[thick] (-0.45,-0.58) -- (0.45,-0.58);
    \draw[thick] (-0.28,-0.82) -- (0.28,-0.82);
    \node[left] at (-0.55,-0.70) {$U_{\mathrm s}$};
    \node[right] at (0.48,-0.58) {$+$};
    \node[right] at (0.31,-0.82) {$-$};
    \node[note,anchor=west] at (0.92,-0.70) {跨过电源：电势升高};

    \draw[current] (0.38,1.58) -- (0.92,1.58);
    \node[red,above] at (0.65,1.58) {$I$};
    \draw[volt] (1.0,0.63) -- (2.0,0.63);
    \node[blue,below] at (1.5,0.58) {$U_1$};
    \draw[volt] (3.1,0.63) -- (4.1,0.63);
    \node[blue,below] at (3.6,0.58) {$U_2$};

    \node[note,align=center] at (2.55,-1.70) {沿电流绕行一周\\[2pt]$U_{\mathrm s}=U_1+U_2$};
  \end{scope}
\end{tikzpicture}


\figcaption{节点电流守恒与回路电压守恒}\label{fig:kirchhoff-circuit-laws}
\end{center}

串联、并联公式不是孤立结论，它们背后分别对应两条守恒思想。

\textbf{节点电流守恒：} 稳恒电路中，导线交汇处不能持续积累电荷。若某节点有若干电流流入、若干电流流出，则
\[
\sum I_{\text{入}}=\sum I_{\text{出}}.
\]
并联电路中“干路电流等于各支路电流之和”，对应节点处的电荷守恒。

\textbf{回路电压守恒：} 一个电荷沿闭合回路走一圈又回到出发点，电势的总变化为零。图中暂用 $U_{\mathrm s}$ 表示理想电源提供的电压；沿约定绕行方向跨过电源时电势升高，沿电流方向通过电阻时电势降低，因此闭合回路中
\[
\sum U_{\text{升}}=\sum U_{\text{降}}.
\]
图示简单回路满足 $U_{\mathrm s}=U_1+U_2$。串联电路中“总电压等于各分电压之和”，对应回路中的能量守恒。实际电源内部的电势变化和电动势将在下一节专门讨论。

\textbf{学习要求：} 高中阶段遇到能直接化简的串并联电路，优先用串并联关系；遇到分叉较多、难以直接化简的电路，可以先标出节点和回路，再用“电流守恒、能量守恒”的思想建立方程。

\end{conceptbox}

\begin{conceptbox}{电功与电功率}

\textbf{电功：} 电荷量 $q$ 通过电势差为 $U$ 的元件时，电场力做功 $W=qU$。恒定电流满足 $q=It$，所以
\[
W=qU=UIt.
\]
这说明电功的表达式来自“每单位电荷转化的能量”与“通过元件的电荷量”的乘积。电功率为
\[
P=\frac{W}{t}=UI.
\]

\textbf{焦耳定律（电热）：}
\[
Q = I^2 Rt
\]

\textbf{焦耳定律的来源：} 对纯电阻，电功全部转化为内能；当导体满足欧姆定律时，$U=IR$，所以发热功率
\[
P_{\text{热}}=UI=I^2R.
\]
若所研究时段内 $I$、$R$ 保持不变，经过时间 $t$ 产生的热量为
\[
Q=P_{\text{热}}t=I^2Rt.
\]
在非纯电阻用电器中，若其线圈或内部电阻可看作恒定，$I^2Rt$ 仍表示这部分产生的电热，但 $UIt$ 还包含转化为机械能、化学能等其他形式的能量。

\textbf{电功率：}
\[
\boxed{P = IU}
\]

\begin{itemize}
  \item 纯电阻电路（如定值电阻、电炉丝；白炽灯丝只在温度近似稳定的小范围内可近似处理）：$P = IU = I^2R = \dfrac{U^2}{R}$。
  \item 在只考虑线圈电阻发热的理想化电动机模型中：输入电功率 $P_{\text{入}}=IU$，线圈发热功率 $P_{\text{热}}=I^2R$，差值 $P_{\text{转}}=UI-I^2R$ 是电磁转换功率。还要忽略铁损、摩擦和风阻等损耗，才能把它等同于轴上的输出机械功率。
\end{itemize}

\end{conceptbox}

\begin{samepage}
\section*{四、公式与适用条件}

\textbf{列式前检查：}
\nopagebreak[4]
\begin{enumerate}
  \item 先画清电流方向和节点，判断元件是串联、并联还是混联。
  \item 使用 $I=U/R$ 前，确认研究对象的端口伏安关系在当前条件下是线性的；“是否满足欧姆定律”和“电能是否全部转化为内能”是两个不同问题。
  \item 非纯电阻用电器中，$P=UI$ 是输入功率，$I^2R$ 只是可等效为电阻的部分所产生的热功率。
  \item 电阻定律中的 $\rho$、$l$、$S$ 描述材料和几何，不代表电路连接关系。
\end{enumerate}
\end{samepage}

\begin{center}
\begin{tabular}{lll}
\toprule
公式 & 物理意义 & 适用条件 \\
\midrule
$I = \dfrac{U}{R}$ & 欧姆定律 & 温度等状态不变的线性导体 \\
$R = \rho\dfrac{l}{S}$ & 电阻定律 & 材料、粗细均匀，温度等状态不变 \\
$W = UIt$ & 电功 & 恒定电流，取元件两端电压和通过电流 \\
$P = IU$ & 电功率 & 元件端口功率，取元件两端电压和通过电流 \\
$Q = I^2Rt$ & 焦耳定律 & 所研究时段内 $I$、$R$ 不变 \\
串：$R = R_1 + R_2$ & 串联总电阻 & 串联 \\
并：$\dfrac{1}{R} = \dfrac{1}{R_1} + \dfrac{1}{R_2}$ & 并联总电阻 & 并联 \\
$\sum I_{\text{入}}=\sum I_{\text{出}}$ & 节点电流关系 & 稳恒电路 \\
$\sum U_{\text{升}}=\sum U_{\text{降}}$ & 回路电压关系 & 闭合回路 \\
\bottomrule
\end{tabular}
\end{center}

\begin{learningbox}{电路题的对象分析清单}
电路题中的 $U$、$I$、$R$、$P$ 都必须带着“对象”使用。列式前先在图上或草稿中写清：
\begin{enumerate}
  \item \textbf{整段电路：} 用 $U_{\text{总}}$、$I_{\text{总}}$、$R_{\text{总}}$，适合先求等效电阻和干路电流。
  \item \textbf{某个支路：} 并联支路常共用同一个电压，但支路电流不同；不要把干路电流直接代入支路电阻。
  \item \textbf{某个元件：} $P=UI$ 中的 $U$、$I$ 必须是同一个元件两端的电压和通过它的电流。
\end{enumerate}
如果对象是电动机、二极管、传感器等元件，先判断整个元件的端口伏安关系能否近似看成线性；不能确认时，优先使用题目给出的伏安关系或图像关系。电动机线圈的电阻 $R$ 可用于计算 $I^2R$ 发热，但不能据此把整个电动机写成 $U=IR$。
\end{learningbox}

\section*{五、典型模型}

\subsection*{模型 1：微观电流与材料电阻}

\textbf{微观电流例：} 某金属中自由电子数密度
$n=8.5\times10^{28}\ \mathrm{m^{-3}}$，导线横截面积
$S=1.0\times10^{-6}\ \mathrm{m^2}$，电子定向移动速率大小
$v_{\mathrm d}=1.0\times10^{-4}\ \mathrm{m/s}$。取
$e=1.6\times10^{-19}\ \mathrm C$，电流大小为
\[
I=neSv_{\mathrm d}
=8.5\times10^{28}\times1.6\times10^{-19}
\times1.0\times10^{-6}\times1.0\times10^{-4}
=1.36\ \mathrm A.
\]
\textbf{关键辨析：} 式中的 $v_{\mathrm d}$ 是电子定向移动速率的大小；约定电流方向与电子定向移动方向相反。

\textbf{电阻变化例：} 同种材料、温度不变时，把均匀导线长度变为原来的 $2$ 倍、直径变为原来的 $\dfrac12$。横截面积变为原来的 $\dfrac14$，所以
\[
\frac{R'}{R}
=\frac{l'}{l}\frac{S}{S'}
=2\times4=8.
\]
\textbf{关键步骤：} 题目改变的是直径时，要先由 $S\propto d^2$ 判断横截面积怎样变化。

\subsection*{模型 2：串联分压 / 并联分流}

串分压：电压与电阻成正比，大电阻分大电压。
并分流：电流与电阻成反比，大电阻分小电流。

\textbf{串联例：} $R_1=4\ \Omega$、$R_2=8\ \Omega$ 串联接在 $12\ \mathrm{V}$ 电源两端，求两电阻分得的电压。

\textbf{解：} 串联电路电流相同，电压按电阻成正比分配：
\[
U_1:U_2=R_1:R_2=4:8=1:2.
\]
总电压为 $12\ \mathrm{V}$，所以
\[
U_1=4\ \mathrm{V},\qquad U_2=8\ \mathrm{V}.
\]
\textbf{并联例：} 若把同样的两个电阻并联接在 $12\ \mathrm V$ 电源两端，则
\[
I_1=\frac{12}{4}=3.0\ \mathrm A,\qquad
I_2=\frac{12}{8}=1.5\ \mathrm A,
\]
所以 $I_1:I_2=2:1=R_2:R_1$，干路电流
$I=I_1+I_2=4.5\ \mathrm A$。

\textbf{提醒：} 串联分压不是“平均分压”，而是按电阻比例分；并联分流也不是平均分流，而是按电阻的反比分配。

\subsection*{模型 3：电路化简与守恒校验}

找等电位点 $\to$ 判断串并联关系 $\to$ 逐级求等效电阻。

\textbf{化简分析清单：}
\begin{enumerate}
  \item 先把导线直接相连的点标成同一个节点，同一节点上的点电势相同。
  \item 再找“只首尾相接、无分叉”的串联部分，和“接在同两个节点之间”的并联部分。
  \item 求出等效电阻后，从总电流或总电压开始，按原电路结构逐级反推各元件的 $I$、$U$、$P$。
\end{enumerate}

\textbf{守恒关系微例：} 某节点有 $I_1=1.2\ \mathrm A$、$I_2=0.8\ \mathrm A$ 流入，只有 $I_3$ 流出，则由电荷守恒
\[
I_3=I_1+I_2=2.0\ \mathrm A.
\]
在另一简单回路中，理想电源电压 $U_{\mathrm s}=12\ \mathrm V$，$R_1=2\ \Omega$、$R_2=4\ \Omega$ 串联，则
\[
I=\frac{U_{\mathrm s}}{R_1+R_2}=2.0\ \mathrm A,\qquad
U_1+U_2=4+8=12\ \mathrm V=U_{\mathrm s}.
\]
节点关系检查电流是否守恒，回路关系检查电势变化是否闭合；它们不仅用于列方程，也可用于复核串并联计算。

\subsection*{模型 4：非纯电阻电路}

在只考虑线圈电阻发热的理想化电动机模型中，
$P_{\text{电}} = UI$，$P_{\text{热}} = I^2R$，电磁转换功率
$P_{\text{转}} = UI - I^2R$。只有题目说明忽略铁损、摩擦、风阻等其他损耗时，才可取 $P_{\text{出}}=P_{\text{转}}$。

\begin{center}
% TikZ 图：非纯电阻电路模型
% Day 19：恒定电流
\begin{tikzpicture}[scale=0.92,>=Latex]
  \tikzset{wire/.style={thick}}
  \draw[wire] (0,0) -- (0,1.4) -- (1.2,1.4);
  \draw[wire] (2.5,1.4) -- (4.2,1.4) -- (4.2,0) -- (0,0);
  \draw[wire] (1.2,1.75) rectangle (2.5,1.05);
  \node at (1.85,1.4) {电动机};
  \draw[wire] (-0.2,0.30) -- (0.2,0.30);
  \draw[wire] (-0.32,0.68) -- (0.32,0.68);
  \node[left] at (-0.45,0.68) {$+$};
  \node[left] at (-0.45,0.30) {$-$};
  \draw[->,red,thick] (2.9,1.75) -- (3.75,1.75) node[midway,above] {$I$};
  \node[font=\small] at (1.85,0.65) {线圈电阻 $R$};

  \node[draw,rounded corners=2pt,align=center,minimum width=1.0cm] (split) at (6.35,1.4) {能量\\分流};
  \draw[->,thick] (4.85,1.4) -- (split.west);
  \node[font=\small,above,anchor=east] at (5.55,2.02) {$P_{\text{电}}=UI$};
  \draw[->,blue,thick] (split.east) -- (8.25,2.18) node[right] {$P_{\text{机}}$};
  \draw[->,orange!85!black,thick] (split.east) -- (8.25,0.62) node[right] {$P_{\text{热}}=I^2R$};
  \node[font=\small] at (6.15,-0.20) {$UI=I^2R+P_{\text{机}}$};
\end{tikzpicture}


\figcaption{非纯电阻电路模型：输入电功率分成发热功率和电磁转换功率}\label{fig:non-ohmic-motor-model}
\end{center}

\textbf{关键判断：} 电动机、电解槽、正在充电的电池都不是纯电阻元件。它们消耗的总电功率是 $UI$，其中只有 $I^2R$ 是线圈或内部电阻发热，不能用 $U^2/R$ 计算总功率。

\subsection*{模型 5：伏安特性曲线与二极管}

欧姆定律适用于线性导体，特征是 $I$-$U$ 图像近似为过原点的直线。二极管是典型的非线性元件，具有明显的\textbf{单向导电性}：正向电压达到一定程度后，电流迅速增大；反向接法时，在高中常见范围内电流很小。

\begin{center}
% TikZ 图：二极管与非线性伏安特性
\begin{tikzpicture}[scale=0.86,>=Latex,line cap=round,line join=round]
  % Left: simple diode circuit.
  \begin{scope}[shift={(-3.2,0)}]
    \draw[thick] (-1.5,1.2) -- (-0.55,1.2);
    \draw[thick] (0.25,1.2) -- (1.5,1.2) -- (1.5,-1.2) -- (-1.5,-1.2) -- (-1.5,-0.05);
    \draw[thick] (-1.5,0.32) -- (-1.5,1.2);
    \draw[thick] (-1.75,0.32) -- (-1.25,0.32);
    \draw[thick] (-1.62,-0.05) -- (-1.38,-0.05);
    \node[font=\small] at (-2.05,0.15) {电源};
    \node[font=\small] at (0.62,1.72) {二极管};
    \draw[thick,fill=gray!15] (-0.55,0.90) -- (-0.55,1.50) -- (0.08,1.20) -- cycle;
    \draw[thick] (0.25,0.90) -- (0.25,1.50);
    \draw[->,red,thick] (-1.35,1.20) -- (-0.72,1.20);
    \node[red,font=\small] at (-1.05,1.53) {正向电流};
    \node[font=\small,align=center] at (0,-1.72) {电路符号：只允许\\一个方向容易通过};
  \end{scope}

  % Right: I-U curve.
  \begin{scope}[shift={(2.45,-0.15)}]
    \draw[->,thick] (-2.35,0) -- (2.95,0) node[right] {$U$};
    \draw[->,thick] (0,-0.65) -- (0,2.65) node[above] {$I$};
    \node[below left] at (0,0) {$O$};

    \draw[blue,very thick] (-2.05,-0.06) .. controls (-1.1,-0.09) and (-0.35,-0.05) .. (0,0);
    \draw[blue,very thick] (0,0) .. controls (0.55,0.04) and (0.90,0.22) .. (1.10,0.70)
      .. controls (1.27,1.18) and (1.43,1.82) .. (1.58,2.35);
    \draw[dashed] (1.05,0) -- (1.05,0.58);
    \node[below] at (1.05,0) {$U_0$};
    \node[blue,font=\small,align=left] at (2.35,1.65) {正向电压较大后\\电流迅速增大};
    \node[font=\small,align=center] at (0,-1.10) {二极管的 $I$-$U$ 图像不是过原点直线};
  \end{scope}
\end{tikzpicture}


\figcaption{二极管的伏安特性：非线性、单向导电，不能直接用 $R=U/I$ 当作恒定电阻}\label{fig:diode-iv-characteristic}
\end{center}

\textbf{学习边界：} 高中阶段不要求解释半导体内部的能带结构，但要会从伏安特性曲线判断元件是否满足欧姆定律。先看图线是否过原点且近似直线；若不是，或电阻随电压、电流明显变化，就不能把它当成恒定电阻处理。

\subsection*{模型 6：传感器的分压读数}

许多传感器不是直接“显示温度或光强”，而是先让某个电学量随外界条件改变。例如热敏电阻、光敏电阻和应变片的电阻会随温度、光照或形变变化。把它们接入分压电路，就能把物理量变化转化为输出电压变化。

\begin{center}
% TikZ 图：传感器分压模型
% Day 19：恒定电流与基本电路
\begin{tikzpicture}[scale=0.92,>=Latex,line cap=round,line join=round]
  \path[use as bounding box] (-4.55,-2.15) rectangle (4.55,2.10);

  \tikzset{
    wire/.style={thick},
    res/.style={thick,fill=white},
    signal/.style={->,very thick,blue},
    current/.style={->,very thick,red}
  }

  % Circuit frame: keep the circuit uncluttered; explanations are placed below.
  \draw[wire] (-3.10,1.30) -- (-1.55,1.30);
  \draw[res] (-1.55,1.00) rectangle (-0.35,1.60);
  \node at (-0.95,1.30) {$R_0$};
  \draw[wire] (-0.35,1.30) -- (1.45,1.30);
  \draw[res] (1.45,0.18) rectangle (2.05,1.30);
  \node at (1.75,0.74) {$R_x$};
  \draw[wire] (1.75,0.18) -- (1.75,-0.55) -- (-3.10,-0.55);
  \draw[wire] (-3.10,-0.55) -- (-3.10,-0.38);
  \draw[wire] (-3.10,0.28) -- (-3.10,1.30);

  % Battery
  \draw[thick] (-3.45,0.28) -- (-2.75,0.28);
  \draw[thick] (-3.30,-0.38) -- (-2.90,-0.38);
  \node[left] at (-3.48,0.28) {$+$};
  \node[left] at (-3.48,-0.38) {$-$};
  \node[left] at (-3.52,-0.05) {$E$};

  % Current and output voltage
  \draw[current] (-2.25,1.52) -- (-1.72,1.52);
  \node[red,font=\small] at (-2.00,1.82) {$I$};
  \draw[signal] (2.65,0.18) -- (2.65,1.30);
  \draw[wire] (1.75,1.30) -- (2.75,1.30);
  \draw[wire] (1.75,0.18) -- (2.75,0.18);
  \node[blue,right] at (2.70,0.74) {$U_{\text{出}}$};

  % Input physical quantities
  \node[font=\small,align=center] at (0,-1.20) {温度、光照、压力等改变 $R_x$};
  \node[font=\small,align=center] at (0,-1.65) {物理量变化 $\to$ 电阻变化 $\to$ 输出电压变化};
\end{tikzpicture}


\figcaption{传感器分压模型：外界物理量改变 $R_x$，电路把变化转化为输出电压}\label{fig:sensor-voltage-divider}
\end{center}

若输出电压取在传感器 $R_x$ 两端，则
\[
U_{\text{出}}=U\frac{R_x}{R_0+R_x}.
\]
因此 $R_x$ 增大时，$U_{\text{出}}$ 增大；$R_x$ 减小时，$U_{\text{出}}$ 减小。若输出电压取在固定电阻 $R_0$ 两端，变化方向则相反。

\textbf{学习边界：} 本节先把传感器看成“可变电阻 + 分压电路”，重点判断外界量变化会让 $R_x$ 增大还是减小，再判断输出电压如何变化。真实传感器还需要标定、放大和数据处理，这些会在实验与科学探究中继续出现。

\section*{六、例题精讲}

\begin{examplebox}{串并联电路的计算}

\textbf{题意分析：} $R_1 = 6\ \Omega$，$R_2 = 12\ \Omega$ 并联后与 $R_3 = 4\ \Omega$ 串联，接在 $U = 12\ \mathrm{V}$ 的电源上。求：
\begin{enumerate}
  \item 电路总电阻。
  \item 干路电流。
  \item $R_1$、$R_2$、$R_3$ 各自的电压和电流。
\end{enumerate}

\textbf{物理模型：} 混联电路。

\begin{center}
% TikZ 图：混联电路例题
% Day 19：恒定电流
\begin{tikzpicture}[scale=0.92,>=Latex]
  \coordinate (A) at (0,0);
  \coordinate (B) at (2.2,0);
  \coordinate (C) at (5.0,0);
  \coordinate (D) at (7.0,0);

  \draw[thick] (A) -- (0,1.2) -- (B |- 0,1.2);
  \draw[thick] (B) -- (B |- 0,1.2);
  \draw[thick] (A) -- (0,-1.2) -- (B |- 0,-1.2);
  \draw[thick] (B) -- (B |- 0,-1.2);

  \draw[thick] (0.55,0.95) rectangle (1.65,1.45);
  \node[above] at (1.10,1.45) {$R_1=6\ \Omega$};
  \draw[thick] (0.55,-1.45) rectangle (1.65,-0.95);
  \node[above] at (1.10,-0.88) {$R_2=12\ \Omega$};

  \draw[thick] (B) -- (C);
  \draw[thick] (3.05,-0.25) rectangle (4.15,0.25);
  \node[above] at (3.60,0.35) {$R_3=4\ \Omega$};
  \draw[thick] (C) -- (D);

  \draw[thick] (D) -- (7,-1.8) -- (-0.25,-1.8);
  \draw[thick] (-0.65,-1.8) -- (-1.2,-1.8) -- (-1.2,0) -- (A);
  \draw[thick] (-0.65,-2.1) -- (-0.65,-1.5);
  \draw[thick] (-0.25,-2.0) -- (-0.25,-1.6);
  \node[above left] at (-0.65,-1.50) {$+$};
  \node[above right] at (-0.25,-1.60) {$-$};
  \node[below] at (-0.45,-2.15) {$12\ \mathrm{V}$};

  \fill (A) circle (1.5pt) node[left] {$a$};
  \fill (B) circle (1.5pt) node[below right] {$b$};
  \node at (1.1,0) {$R_{12}$};
  \draw[->,red,thick] (5.8,0.35) -- (6.6,0.35) node[midway,above] {$I$};
\end{tikzpicture}


\figcaption{混联电路：先合并并联部分，再与 $R_3$ 串联}\label{fig:mixed-circuit-example}
\end{center}

\textbf{解题过程：}

\begin{enumerate}
  \item $R_{12} = \dfrac{R_1 R_2}{R_1 + R_2} = \dfrac{6 \times 12}{6 + 12} = 4\ \Omega$

  $R_{\text{总}} = R_{12} + R_3 = 4 + 4 = 8\ \Omega$

  \item $I_{\text{总}} = \dfrac{U}{R_{\text{总}}} = \dfrac{12}{8} = 1.5\ \mathrm{A}$

  \item $U_3 = I_{\text{总}} R_3 = 1.5 \times 4 = 6\ \mathrm{V}$

  $U_{12} = U - U_3 = 12 - 6 = 6\ \mathrm{V}$（并联部分电压）

  $I_1 = \dfrac{U_{12}}{R_1} = \dfrac{6}{6} = 1\ \mathrm{A}$，$I_2 = \dfrac{U_{12}}{R_2} = \dfrac{6}{12} = 0.5\ \mathrm{A}$
\end{enumerate}

\textbf{方法总结：} 混联电路的解题顺序：先合并并联部分 $\to$ 求总电阻 $\to$ 求总电流 $\to$ 逐级反推各支路。

\end{examplebox}

\begin{examplebox}{电功率计算}

\textbf{题意分析：} 标有“$220\ \mathrm{V}$、$100\ \mathrm{W}$”的灯泡：
\begin{enumerate}
  \item 正常发光时电流和电阻是多少？
  \item 接在 $110\ \mathrm{V}$ 电源上时实际功率是多少？（假设电阻不变）
\end{enumerate}

\textbf{物理模型：} 纯电阻电路。

\textbf{解题过程：}

\begin{enumerate}
  \item $I = \dfrac{P}{U} = \dfrac{100}{220} \approx 0.455\ \mathrm{A}$

  $R = \dfrac{U^2}{P} = \dfrac{220^2}{100} = 484\ \Omega$

  \item $P_{\text{实}} = \dfrac{U_{\text{实}}^2}{R} = \dfrac{110^2}{484} = \dfrac{12100}{484} = 25\ \mathrm{W}$
\end{enumerate}

\textbf{方法总结：} 功率题先看题目是否允许电阻不变；若说明“假设电阻不变”，可用同一个 $R$，实际灯丝则会随温度变化。

\end{examplebox}

\begin{examplebox}{非纯电阻电路}

\textbf{题意分析：} 一台电动机的线圈电阻 $R = 1\ \Omega$，接在 $U = 220\ \mathrm{V}$ 电源上工作时电流 $I = 10\ \mathrm{A}$。求：
\begin{enumerate}
  \item 输入电功率。
  \item 发热功率。
  \item 电磁转换功率和效率。（忽略铁损、摩擦、风阻等其他损耗。）
\end{enumerate}

\textbf{物理模型：} 非纯电阻电路，$P = UI \neq I^2R$。

\begin{center}
% TikZ 图：非纯电阻电路例题
% Day 19：恒定电流
\begin{tikzpicture}[scale=0.92,>=Latex]
  \tikzset{wire/.style={thick}}
  \draw[wire] (0,0) -- (0,0.30);
  \draw[wire] (0,0.68) -- (0,1.35) -- (1.4,1.35);
  \draw[wire] (3.0,1.35) -- (4.5,1.35) -- (4.5,0) -- (0,0);
  \draw[wire] (1.4,1.75) rectangle (3.0,0.95);
  \node at (2.2,1.35) {电动机};
  \node[below] at (2.2,0.95) {$R=1\ \Omega$};
  \draw[wire] (-0.2,0.30) -- (0.2,0.30);
  \draw[wire] (-0.32,0.68) -- (0.32,0.68);
  \node[left] at (-0.48,0.68) {$+$};
  \node[left] at (-0.48,0.30) {$-$};
  \node[left] at (-0.55,0.48) {$220\ \mathrm{V}$};
  \draw[->,red,thick] (3.25,1.70) -- (4.15,1.70) node[midway,above] {$I=10\ \mathrm{A}$};

  \node[draw,rounded corners=2pt,align=left,font=\small] at (7.0,0.80)
    {$P_{\text{电}}=UI$\\$P_{\text{热}}=I^2R$\\$P_{\text{转}}=UI-I^2R$};
\end{tikzpicture}


\figcaption{非纯电阻例题图：电动机输入功率不等于线圈发热功率}\label{fig:non-ohmic-motor-example}
\end{center}

\textbf{解题过程：}

\begin{enumerate}
  \item $P_{\text{电}} = UI = 220 \times 10 = 2200\ \mathrm{W}$

  \item $P_{\text{热}} = I^2 R = 10^2 \times 1 = 100\ \mathrm{W}$

  \item $P_{\text{转}} = P_{\text{电}} - P_{\text{热}} = 2200 - 100 = 2100\ \mathrm{W}$。题设忽略其他损耗，故输出机械功率也为 $2100\ \mathrm{W}$。

  因此 $P_{\text{出}}=P_{\text{转}}=2100\ \mathrm{W}$，效率
  \[
  \eta = \frac{P_{\text{出}}}{P_{\text{电}}} = \frac{2100}{2200} \approx 95.5\%.
  \]
\end{enumerate}

\textbf{方法总结：} 非纯电阻题先分清输入功率、发热功率和电磁转换功率：输入功率用 $P = UI$，线圈发热功率用 $P = I^2R$；是否能把差值当作轴输出，还要看题目是否忽略铁损和机械损耗。

\end{examplebox}

\section*{七、变式训练}

\begin{enumerate}
  \item 某金属导线中自由电子数密度
  $n=5.0\times10^{28}\ \mathrm{m^{-3}}$，横截面积
  $S=2.0\times10^{-6}\ \mathrm{m^2}$，电子定向移动速率大小
  $v_{\mathrm d}=1.25\times10^{-4}\ \mathrm{m/s}$。取
  $e=1.6\times10^{-19}\ \mathrm C$，求电流大小，并说明电子定向移动方向与约定电流方向的关系。

  \item 同种材料、温度不变时，一段均匀导线的长度变为原来的 $3$ 倍、直径变为原来的 $\dfrac12$，电阻变为原来的多少倍？

  \item $R_1 = 3\ \Omega$、$R_2 = 6\ \Omega$ 串联后接在 $U = 18\ \mathrm{V}$ 的电源上，求总电阻、电流和 $R_1$、$R_2$ 两端的电压。

  \item 将上题改为并联，求总电阻和干路电流。
  \begin{center}
  \begin{tikzpicture}[scale=0.86,>=Stealth,line cap=round,line join=round]
  \begin{scope}
    \draw[thick] (0,0) -- (0,1.2) -- (0.75,1.2);
    \draw[thick] (1.35,1.2) -- (2.0,1.2);
    \draw[thick] (2.6,1.2) -- (3.35,1.2) -- (3.35,0) -- (1.28,0);
    \draw[thick] (0.78,0) -- (0,0);
    \draw[thick] (0.75,0.92) rectangle (1.35,1.48);
    \draw[thick] (2.0,0.92) rectangle (2.6,1.48);
    \node[font=\small] at (1.05,1.2) {$R_1$};
    \node[font=\small] at (2.3,1.2) {$R_2$};
    \draw[thick] (0.90,-0.16) -- (0.90,0.16);
    \draw[thick] (1.14,-0.26) -- (1.14,0.26);
    \node[font=\small] at (1.02,-0.45) {$18\,\mathrm V$};
    \node[font=\small] at (1.65,1.82) {串联：电流相同};
  \end{scope}
  \begin{scope}[xshift=4.6cm]
    \draw[thick] (0,0) -- (0,1.6) -- (0.65,1.6);
    \draw[thick] (1.35,1.6) -- (2.2,1.6) -- (2.2,0) -- (1.12,0);
    \draw[thick] (0.62,0) -- (0,0);
    \draw[thick] (0,0.6) -- (0.65,0.6);
    \draw[thick] (1.35,0.6) -- (2.2,0.6);
    \draw[thick] (0.65,1.32) rectangle (1.35,1.88);
    \draw[thick] (0.65,0.32) rectangle (1.35,0.88);
    \node[font=\small] at (1.0,1.6) {$R_1$};
    \node[font=\small] at (1.0,0.6) {$R_2$};
    \draw[thick] (0.74,-0.16) -- (0.74,0.16);
    \draw[thick] (0.98,-0.26) -- (0.98,0.26);
    \node[font=\small] at (0.88,-0.45) {$18\,\mathrm V$};
    \node[font=\small] at (1.1,2.22) {并联：电压相同};
  \end{scope}
\end{tikzpicture}

  \end{center}

  \item 一台电动机线圈电阻 $R = 2\ \Omega$，工作电压
  $U = 110\ \mathrm{V}$，电流 $I = 20\ \mathrm{A}$。若忽略铁损、摩擦和风阻等其他损耗，求输入功率、线圈发热功率、输出机械功率和效率。
  \begin{center}
  \begin{tikzpicture}[scale=0.86,>=Stealth,line cap=round,line join=round]
  \path[use as bounding box] (-0.15,-0.85) rectangle (7.65,2.15);
  \draw[thick] (0,0) -- (0,1.2) -- (0.9,1.2);
  \draw[thick] (2.1,1.2) -- (3.2,1.2) -- (3.2,0) -- (1.2,0);
  \draw[thick] (0.7,0) -- (0,0);
  \draw[thick] (1.5,1.2) circle (0.6);
  \node at (1.5,1.2) {$M$};
  \draw[->,blue,thick] (0.32,1.48) -- (0.92,1.48) node[midway,above] {$I$};
  \draw[thick] (0.82,-0.26) -- (0.82,0.26);
  \draw[thick] (1.06,-0.16) -- (1.06,0.16);
  \node[font=\small] at (0.94,-0.48) {$U$};
  \node[font=\small] at (1.55,-0.55) {电动机电路};
  \draw[->,red,thick] (4.0,1.2) -- (4.95,1.2) node[midway,above] {$UI$ 输入};
  \draw[fill=gray!10,thick] (5.05,0.85) rectangle (5.45,1.55);
  \draw[->,orange,thick] (5.45,1.35) -- (6.75,1.95)
    node[right,font=\small] {机械输出};
  \draw[->,orange,thick] (5.45,1.05) -- (6.75,0.45)
    node[right,font=\small] {发热 $I^2R$};
  \node[font=\scriptsize] at (5.85,-0.18) {非纯电阻：};
  \node[font=\scriptsize] at (5.85,-0.78) {$UI$ 不是全部变成热};
\end{tikzpicture}

  \end{center}

  \item 某元件的 $I$-$U$ 图像不是过原点直线，能否把 $R=U/I$ 当作该元件的恒定电阻？为什么？

  \item 某传感器接成分压电路，输出电压取在传感器两端。若外界条件使传感器电阻增大，输出电压怎样变化？

  \item 某节点有 $0.6\ \mathrm A$、$0.9\ \mathrm A$ 两股电流流入，有 $0.4\ \mathrm A$ 和 $I_x$ 两股电流流出，求 $I_x$。另有一简单闭合回路，理想电源电压为 $9\ \mathrm V$，两个电阻的电压降分别为 $3\ \mathrm V$ 和 $U_2$，求 $U_2$。
\end{enumerate}

\begin{learningbox}{变式训练参考结果}
\begin{enumerate}[itemsep=0.45em,topsep=0.3em,parsep=0pt]
  \item $I=neSv_{\mathrm d}=2.0\ \mathrm A$；电子定向移动方向与约定电流方向相反。
  \item 横截面积变为原来的 $1/4$，所以
  $R'/R=(l'/l)(S/S')=3\times4=12$。
  \item $R_{\text{总}}=9\ \Omega$，$I=2.0\ \mathrm A$，
  $U_1=6\ \mathrm V$，$U_2=12\ \mathrm V$。
  \item $R_{\text{总}}=R_1R_2/(R_1+R_2)=2.0\ \Omega$，干路电流
  $I=9.0\ \mathrm A$。
  \item 输入功率 $P_{\text{入}}=UI=2.20\times10^3\ \mathrm W$，发热功率
  $P_{\text{热}}=I^2R=800\ \mathrm W$；由题设忽略其他损耗，输出机械功率为 $1.40\times10^3\ \mathrm W$，
  效率约为 $63.6\%$。
  \item 不能。$R=U/I$ 只能表示该工作点的电阻值；图线不是过原点直线，说明
  $U/I$ 随工作状态变化，元件不是定值电阻。
  \item 设与传感器串联的定值电阻为 $R_0$，则
  $U_{\text{出}}=UR_{\text{传}}/(R_0+R_{\text{传}})$，所以传感器电阻增大时输出电压增大。
  \item 节点电流守恒给出 $0.6+0.9=0.4+I_x$，故
  $I_x=1.1\ \mathrm A$；回路电压守恒给出
  $9=3+U_2$，故 $U_2=6\ \mathrm V$。
\end{enumerate}
\end{learningbox}

\section*{八、竞赛入门补充}

\begin{competitionbox}
\textbf{电路的化简方法——对称法}

如果电路具有对称性，可以利用对称性简化计算。

\textbf{例题：} 一个正方体框架每条棱的电阻为 $R$，求体对角线两个顶点间的等效电阻。

\begin{center}
% TikZ 图：正方体电阻网络
% Day 19：恒定电流
\begin{tikzpicture}[scale=0.86,>=Latex]
  \coordinate (A) at (0,0);
  \coordinate (B) at (2.2,0);
  \coordinate (C) at (2.2,2.0);
  \coordinate (D) at (0,2.0);
  \coordinate (E) at (0.85,0.75);
  \coordinate (F) at (3.05,0.75);
  \coordinate (G) at (3.05,2.75);
  \coordinate (H) at (0.85,2.75);

  \foreach \p/\q in {A/B,B/C,C/D,D/A,E/F,F/G,G/H,H/E,A/E,B/F,C/G,D/H} {
    \draw[thick] (\p) -- (\q);
  }
  \foreach \p in {A,B,C,D,E,F,G,H} {\filldraw[black] (\p) circle (1.5pt);}
  \node[below left] at (A) {$A$};
  \node[above right] at (G) {$B$};
  \draw[->,red,thick] (-0.9,-0.25) -- (A) node[midway,below] {$I$};
  \draw[->,red,thick] (G) -- (3.95,3.05) node[midway,above] {$I$};
  \node[font=\small,align=center] at (1.65,-0.85) {每条棱电阻均为 $R$，求体对角线 $AB$ 间等效电阻};
\end{tikzpicture}


\figcaption{正方体电阻网络：从体对角线两端接入电流}\label{fig:cube-resistor-network}
\end{center}

\begin{center}
% TikZ 图：正方体电阻网络的等效分层
% Day 19：恒定电流
\begin{tikzpicture}[scale=0.92,>=Latex]
  \node[circle,draw,inner sep=2pt] (A) at (0,0) {$A$};
  \node[circle,draw,inner sep=2pt] (L1) at (2.1,0) {三点};
  \node[circle,draw,inner sep=2pt] (L2) at (4.6,0) {三点};
  \node[circle,draw,inner sep=2pt] (B) at (6.7,0) {$B$};

  \draw[thick] (A) -- (L1) node[midway,above] {$R/3$};
  \draw[thick] (L1) -- (L2) node[midway,above] {$R/6$};
  \draw[thick] (L2) -- (B) node[midway,above] {$R/3$};
  \draw[->,red,thick] (-0.7,0) -- (A);
  \draw[->,red,thick] (B) -- (7.4,0);

  \node[align=center] at (3.35,-1.35)
    {对称点电势相同，可以合并\\三段串联：$R_{\text{总}}=\dfrac{R}{3}+\dfrac{R}{6}+\dfrac{R}{3}$};
\end{tikzpicture}


\figcaption{正方体电阻网络的对称化简：同层等电位点可以合并}\label{fig:cube-equivalent-network}
\end{center}

\textbf{解：} 设电流从体对角线一端 $A$ 流入，从另一端 $B$ 流出。与 $A$ 相邻的三个顶点等电位，可合并为一组；与 $B$ 相邻的三个顶点也等电位，可合并为一组。于是电路化为三段串联：
\[
R_{\text{总}} = \frac{R}{3}+\frac{R}{6}+\frac{R}{3}=\frac{5}{6}R
\]

\textbf{竞赛提示：} 遇到复杂网络电阻问题，先观察几何结构和接入端是否共同保持某种对称性。由对称性可知，同一层的对应顶点电势相同、对应支路电流相等。用理想导线连接这些等势点时，连接线上没有电流，原电路各支路的电流分布不变，因此可以把等势点合并后再化简。不能只因图形外观对称就合并节点，还要检查电源或电流的接入方式是否保留了这种对称性。
\end{competitionbox}

\section*{九、易错点}

\begin{mistakebox}
\begin{enumerate}
  \item \textbf{混联电路中串并联关系判断错：} 先标节点和等电位点，再判断连接关系。两个元件首尾相接且中间节点没有其他支路时，它们串联；两个元件分别接在同一对节点之间时，它们并联。不能只凭图形横着画还是竖着画来判断。

  \item \textbf{把“是否线性”和“是否纯电阻”混为一谈：} 判断能否用 $I=U/R$，要看元件在当前条件下是否满足欧姆定律；判断电能是否全部转化为内能，才看它是否为纯电阻。满足欧姆定律的纯电阻可用 $P=UI=I^2R=U^2/R$，电动机输入功率用 $P=UI$。

  \item \textbf{并联等效电阻算得过大：} 化简后先把结果与每个分电阻比较；只要各支路都是正电阻，就必须有 $R_{\text{总}}<R_{\min}$。$n$ 个相同电阻 $R$ 并联时，再用 $R_{\text{总}}=R/n$ 作快速复核。

  \item \textbf{串联分压时把电压和电流混淆：} 先在各元件旁标同一个电流 $I$，再写 $U_i=IR_i$；因此大电阻分得大电压，而不是大电流。若题图出现分叉，则先别套这一结论，要重新判断是否仍为串联。

  \item \textbf{误以为焦耳定律 $Q = I^2Rt$ 只适用于纯电阻电路：} 当 $I$、$R$ 在所研究时段内不变时，它可用于计算任何电阻性部分产生的电热；但在非纯电阻用电器中，消耗的电能不全变成热，通常 $Q \neq IUt$。若 $I$ 或 $R$ 随时间变化，则要分段求和，不能直接使用一个不变的 $I^2R$ 乘总时间。

  \item \textbf{把二极管当普通电阻：} 二极管的伏安特性是非线性的，且正反向性质不同，不能直接套串并联电阻公式。

  \item \textbf{传感器题不看输出端接在哪里：} 先圈出输出电压取在哪两个点之间；输出取在可变电阻两端或固定电阻两端，变化方向可能相反。

  \item \textbf{复杂电路只凭直觉猜串并联：} 分叉较多时先标节点，再用节点电流守恒和回路电压关系检查，最后再决定能否化简。
\end{enumerate}
\end{mistakebox}

\section*{十、今日小结}

\begin{enumerate}
  \item 电流的微观表达式 $I=n\lvert q\rvert Sv_{\mathrm d}$；约定电流方向与金属中电子定向移动方向相反。
  \item 欧姆定律 $I = \dfrac{U}{R}$ 是有条件的实验定律；均匀导体的电阻定律为 $R = \rho\dfrac{l}{S}$。
  \item 串联：$I$ 相等，$U$ 分压，$R = R_1 + R_2$。
  \item 并联：$U$ 相等，$I$ 分流，$\dfrac{1}{R} = \dfrac{1}{R_1} + \dfrac{1}{R_2}$。
  \item 节点电流守恒来自电荷守恒，回路电压守恒来自能量守恒。
  \item 由 $W=qU$、$q=It$ 得 $W=UIt$ 和 $P=UI$；$I$、$R$ 不变时，焦耳定律为 $Q = I^2Rt$。
  \item 非纯电阻电路中 $P_{\text{电}} = UI$，$P_{\text{热}} = I^2R$，两者不等。
  \item 二极管等非线性元件的伏安特性不是直线，不能当作恒定电阻处理。
  \item 热敏、光敏等传感器常可看成可变电阻，借助分压电路把物理量变化转化为电压信号。
\end{enumerate}

\textbf{一分钟复盘：} 任取一个电路，先用同色标出等势节点，再化简串并联关系；接着在每个元件旁标 $I$、$U$、$P$。若电路不能直接化简，就写一条节点电流关系和一条回路电压关系；遇到电动机、二极管或传感器，先分别判断其端口伏安关系是否近似线性、输入电能转化成了哪些形式。
